\section{Dependency and source conventions}
\label{app:v97-dependencies}
The unchanged physical formulation is \eqref{eq:v76-event}, with the
reference \eqref{eq:v76-reference} and pointwise goal
\eqref{eq:v76-target}. The complete preceding article is included in
the technical archive. All 231 earlier core files retain their bytes.

The new collar lemma uses only the section momentum bound and the
radial comparison-flow identity. The capped anchor uses the comparison
coarea density, fixed-time injectivity, and the original guard modulus.
The running-compression theorem uses the original stopped lifetimes,
radial growth, the derivative of both angular-density factors, and the
positive-cutoff source estimate \eqref{eq:v94-cutoff}. That last input
is proved in retained modules 196 and 222 from the finite-count
critical-volume layer cake; its leading constant is cutoff independent,
but its finite-count remainder is not. The cutoff in the new proof is
positive and fixed before every collision limit.

The final recovery theorem combines these two distinct comparisons
and uses the original scalar average and positive remainder. The
support mass estimate additionally uses the original backward-end and
critical-volume estimates. No new theorem on a spectral insertion
outside the exact event, no uniform untruncated first inverse gradient,
and no interchange of a zero-gradient limit with a collision limsup
is required or asserted.

The Borel and signed-measure foundations remain those of the retained
geometric chain, with the standard bounded-variation and definable
integration references \cite{AFP,CMInt,Coste}. Source preservation and
finite diagnostics do not supply an independent reconstruction of
that chain. The new proofs provide an estimated recovered source;
they do not establish vanishing of the remaining critical/short-source
essential height. The two distinct revision-88 variants retain their
separate source identities throughout the archive.
